\documentclass[11pt]{article}
\usepackage[a4paper,margin=25mm]{geometry}
\usepackage{amsmath,amssymb,amsthm,booktabs,longtable,hyperref}
\newtheorem{theorem}{Theorem}
\newtheorem{lemma}{Lemma}
\newcommand{\Z}{\mathbb Z}
\newcommand{\Q}{\mathbb Q}
\newcommand{\vp}{v_p}
\newcommand{\dd}{\mathrm d}
\newcommand{\z}{\zeta(3)}
\DeclareMathOperator{\lcm}{lcm}
\title{A CMF proof of the linear independence of $1,\zeta(3),\zeta(4)$}
\author{Proof development and exact certificates}
\date{6 October 2026}
\begin{document}
\maketitle

\noindent\textbf{Main conclusion.} For every $(A,B)\in\Q^2\setminus\{(0,0)\}$,
$A\zeta(3)+B\zeta(4)$ is irrational. The proof uses the trajectory in the
handoff, with no change of initial point, direction, or product convention.
For the nonvanishing argument we use all three columns of its $3\times3$
block; the displayed scalar candidate remains column zero.
The new ingredient relative to the $\zeta(3)$ argument is a
\emph{second-derivative dual residue functional}. It supplies the missing
arithmetic estimate for the $\zeta(4)$ readout.

\medskip
\noindent\textbf{Literature.} Li Lai's preprint of 5 October 2026 states the
same linear-independence conclusion as Theorem 1.6:
\url{https://arxiv.org/abs/2610.06838}. The argument below is a CMF argument
using the trajectory specified here. No novelty claim is made for the
linear-independence statement itself. This document is a mathematical proof
with exact arithmetic certificates; it has not yet been formalized in Lean.

\section*{I. The precise rational sequences}
Use $\mathrm{MeijerG}(5,2,5,5,1)$ with
\[
 s=(1,1,0,0,2;0,0,0,0,0),\qquad
 v=(11,12,-16,-17,10;0,1,2,-1,-2).
\]
Write
\[
 W_N=M(0)M(1)\cdots M(N-1)
 =\begin{pmatrix}V_N&*\\0&(-1)^N\end{pmatrix},\qquad T_N=V_N^{-T}.
\]
The exact matrix reconstruction and the block identity are certified in
\texttt{verify\_matrix.py}. Put
\[
 U_{34}=\begin{pmatrix}0&-3&3\\2&-2&0\\0&0&6\end{pmatrix},\qquad
 \mathcal R_N=U_{34}T_N
 =\begin{pmatrix}p^{(3)}_{N0}&p^{(3)}_{N1}&p^{(3)}_{N2}\\
 p^{(4)}_{N0}&p^{(4)}_{N1}&p^{(4)}_{N2}\\q_{N0}&q_{N1}&q_{N2}\end{pmatrix}.
\]
These entries are rational; they are not yet integer numerators and denominators.
For column zero, with $t=W_N^{-T}e_0$, their ratios are exactly
\[
 \frac{p^{(3)}_{N0}}{q_{N0}}=\frac{t_2-t_1}{2t_2},\qquad
 \frac{p^{(4)}_{N0}}{q_{N0}}=\frac{t_0-t_1}{3t_2}.
\]
Thus the numerator row for $A\zeta(3)+B\zeta(4)$ is precisely
$(2B,-3A-2B,3A,0)$ and the denominator row is $(0,0,6,0)$.
Also,
\begin{equation}\label{new:det}
 \det\mathcal R_N=36/\det V_N\ne0.
\end{equation}

\section*{II. The analytic input, with the exact readouts}
The technical proof reproduced in Part~VII below establishes the constants
\[
 a=53.28758332262020526735\ldots,\qquad
 b=59.51995592483447507214\ldots
\]
and the estimates
\begin{align}
 |q_{Nj}|&\le e^{aN+o(N)},\label{new:growth}\\
 |q_{Nj}\zeta(3)-p^{(3)}_{Nj}|+
 |q_{Nj}\zeta(4)-p^{(4)}_{Nj}|&\le e^{-bN+o(N)}\quad(j=0,1,2).
 \label{new:errors}
\end{align}
Here is explicitly why its argument gives both errors. Let
$\mathbf A_N,\mathbf B_N,\mathbf C_N$ be the $\zeta(4)$, $\zeta(3)$, and
rational coefficient rows of the original rational kernel, and set
$\mathbf S_N=\zeta(4)\mathbf A_N+\zeta(3)\mathbf B_N+\mathbf C_N$.
The coefficientwise identity is
\[
 \begin{pmatrix}\mathbf A_N\\\mathbf B_N\\\mathbf C_N\end{pmatrix}
 =(-1)^N\begin{pmatrix}3&0&0\\-2&-2&0\\1&1&1\end{pmatrix}V_N.
\]
At $N=0$, direct cross products give
\[
 \mathbf A_0\times\mathbf S_0
   =-(0,3,6\zeta(3)-3),\qquad
 \mathbf B_0\times\mathbf S_0=(-2,2,6\zeta(4)).
\]
Consequently, as row-vector identities,
\[
 \zeta(3)\mathbf q_N-\mathbf p^{(3)}_N
  =-\frac{\mathbf A_N\times\mathbf S_N}{\det V_N},\qquad
 \zeta(4)\mathbf q_N-\mathbf p^{(4)}_N
  =\frac{\mathbf B_N\times\mathbf S_N}{\det V_N}.
\]
The bounds $\|\mathbf A_N\|,\|\mathbf B_N\|\le e^{bN+o(N)}$,
$\|\mathbf S_N\|\le e^{-aN+o(N)}$, and
$\det V_N=e^{(2b-a)N+o(N)}$ prove \eqref{new:errors}.
The small-solution identification, including the complex conjugate modes,
is proved in Part~VII. No lower bound on any oscillating scalar error is used.

\section*{III. The second-derivative identity}
The original proof uses the formal functional $\mathcal L_K$ defined in
Part~VII. For a proper rational function with integer poles at most $K$,
\[
 \mathcal L_K[f]=-\sum_{m>K} f'(m).
\]
For polynomial parts, $\mathcal L_K[t^d]=B_d(K+1)$; the zeta values in
partial fractions are treated as independent formal symbols $Z_r$.
Define a second functional by
\[
 \mathcal M_K[f]:=-\mathcal L_K[f'].
\]
For proper $f$, this is $\sum_{m>K}f''(m)$, and explicitly
\[
 \mathcal M_K[(t-h)^{-r}]
 =r(r+1)\bigl(Z_{r+2}-H_{K-h}^{(r+2)}\bigr),\qquad
 \mathcal M_K[t^d]=-dB_{d-1}(K+1).
\]
It satisfies the exact identities
\begin{equation}\label{new:shift}
 \mathcal M_K[f(t)-f(t-1)]=-f''(K),\qquad
 \mathcal M_K[f]-\mathcal M_{K-1}[f]=-f''(K).
\end{equation}
They follow directly by differentiating the argument of the already proved
$\mathcal L_K$ identities. In particular,
$\mathcal M_K[1]=0$, $\mathcal M_K[t]=-1$, and
$\mathcal M_K[t^2]=-(2K+1)$ are integers for integer $K$.

For the dual parameters $a'=1-a$, $b'=-b$, the initial rational kernel is
$r_0(s)=(s-1)/s^2$. Its four second-functional moments are exactly
\begin{equation}\label{new:g0}
 g^{[2]}_0=(2Z_3-6Z_4,\ 2Z_3,\ -1,\ 0).
\end{equation}
The last two entries are polynomial contributions, not divergent sums.
The same differential-equation telescoper as in Part~VII is
\[
 D(-s)r(s)=g(s)-g(s-1),\qquad
 g(s)=\prod_j(s+b'_j)r(s).
\]
Applying \eqref{new:shift} gives the companion relation provided $g''(K)=0$.
This extra boundary condition must be checked; it does not follow just from
the first-derivative argument.

\begin{lemma}\label{new:transport}
For every $N\ge1$ the dual moments defined by $\mathcal M$ have the exact
same CMF transport as the first-functional moments, with final cutoff $K=2N$.
\end{lemma}
\begin{proof}
Start at $a'_0=(0,0,1,1,-1)$, $b'_0=0$, $K=0$. First decrease $a'_0$ by
$11N$, $a'_1$ by $12N$, and $a'_4$ by $10N$. Initially $s^5r_0(s)$ has a
triple zero at zero, and these steps preserve it, so $g''(0)=0$ throughout.
After these steps, the kernel has triple zeros at every integer
$1,\ldots,2N$: at $1$ two new zeros join the initial factor $s-1$, and at
$2,\ldots,2N$ all three newly inserted blocks vanish. Thus the cutoff may
be changed from $0$ to $2N$ by \eqref{new:shift}, with no boundary terms.

Next increase $a'_2$ by $16N$ and $a'_3$ by $17N$, increase $b'_3$ by $N$
and $b'_4$ by $2N$, then decrease $b'_1$ by $N$ and $b'_2$ by $2N$.
The positive shifts do not remove any zero at $2N$; the first negative lower
block also does not. The last block can remove one factor $s-2N$, leaving
a double zero of $r$, but the corresponding factor in $\prod(s+b'_j)$
restores the third zero of $g$. Thus $g''(2N)=0$ at every vertex.
All poles of $g$ are below the cutoff. The companion coefficient remains
positive, following $4\to4+33N\to4\to4+3N\to4$, and the same explicit
parameter-separation checks as in Part~VII show that every axis step is
invertible. The gamma shift identities therefore transport all four formal
moments exactly. This proves the assertion for every $N$; $N=0$ is
\eqref{new:g0}.
\end{proof}

Let $J_N$ be the integer polynomial concomitant matrix displayed in Part~VII.
Its exact, generically certified transport identity applies to these moments
as well. Direct multiplication at the initial point gives
\[
 g^{[2]}_0J_0[0\!:\!3,:]^T=(-2,2,6Z_4).
\]
The $Z_3$ terms cancel identically. With
\[
 R_N^*(t)=\frac{(t+N+1)_{11N}(t+1)_{11N}(t+2N)_{8N+1}}
 {(t+13N+1)_{15N+1}(t+14N+1)_{15N+1}},\quad
 Q_{Nj}(t)=\sum_{r=0}^3(J_N)_{jr}(-t-12N-1)^r,
\]
we therefore obtain, coefficientwise for all $N$,
\begin{equation}\label{new:dual4}
 (-1)^N\mathcal M_{-10N-1}[Q_{Nj}(t)R_N^*(t)]
      =Z_4q_{Nj}-p^{(4)}_{Nj}.
\end{equation}
The cutoff is $-10N-1$ because $s=t+12N+1$ and the unshifted cutoff is
$2N$. The largest harmonic index in \eqref{new:dual4} is consequently
$19N$, rather than the $18N$ in the first-functional formula.

The polynomial $Q_{Nj}$ has integer coefficients and degree at most three.
Since $R_N^*$ has a monic integer denominator and degree $-1$, the polynomial
part of $Q_{Nj}R_N^*$ has integer coefficients and degree at most two.
Its $\mathcal M$ value is an integer by the formulas above. Thus no unaccounted
polynomial or regularization term occurs in \eqref{new:dual4}.
Exact rational checks of all three columns at $N=0,1,2,4,8,12,20$ accompany
the proof; an additional exact audit verifies every one of the $72+144$ elementary
steps at $N=1,2$, including the triple boundary zeros. The lemma and pairing
identity, rather than these checks, prove
\eqref{new:dual4} for all $N$.

\section*{IV. The all-prime denominator bound}
Use $f_p,\ell_p,V_o,V_d,m_o,m_d,A_p,C_p,F_N,E_o,E_d$ as explicitly defined
in Part~VII, Section~3. These use full Legendre sums at every prime, so the
following definition provides integer multipliers for every $N$.
Define the additional original-row bound
\[
 B_p=\min_{\substack{0\le k\le12N\\m_o(k)\ge2}}
 \{V_o(k)-(m_o(k)-2)\ell_p(k+17N)\},\qquad
 E_{o4}(N,p)=\max\{0,v_p(F_N)-B_p-C_p\}\quad(p>3).
\]
The row with coefficients of $\zeta(3)$ in the original kernel has valuation
at least $B_p$. The cofactor expressing $(1,-1,0)T_N$ uses this row and the
rational row, so $E_{o4}$ bounds the $\zeta(4)$ numerator readout.
For the dual formula put
\[
 h_4(k)=k-10N-1,\qquad L(k)=\max\{k-1,29N+1-k\},
\]
\[
 E_{d4}(N,p)=\max\left\{0,
 \max_{13N+1\le k\le29N+1}
 \bigl((m_d(k)-1)\ell_p(L(k))+3\ell_p(h_4(k))-V_d(k)\bigr)\right\}.
\]
Indeed a partial-fraction coefficient at a pole of order $m$ has valuation
at least $V_d-(m-r)\ell_p(L)$, and the harmonic factor supplied by
$\mathcal M$ loses at most $(r+2)\ell_p(h_4)$.
Since $L\ge h_4$, their sum is bounded below by
$V_d-(m-1)\ell_p(L)-3\ell_p(h_4)$.
The factors $r(r+1)$ are integers and only improve the bound.
The polynomial part is integral, as proved above. This establishes the dual
bound at all primes, including $2$ and $3$.

For $p>3$, set
\[
 E_{34}(N,p)=\max\{\min(E_o,E_d),\ \min(E_{o4},E_{d4})\},
\]
and for $p=2,3$ set $E_{34}(N,p)=\max\{E_d,E_{d4}\}$.
Finally put
\begin{equation}\label{new:deltaN}
 \Delta_N=\prod_{p\le29N+1}p^{E_{34}(N,p)}\quad(N\ge1),\qquad\Delta_0=1.
\end{equation}
The first pair of bounds clears $p^{(3)}_{Nj}$ and $q_{Nj}$ (they are three
times the corresponding previous readouts), and the second pair clears
$p^{(4)}_{Nj}$. Therefore
\begin{equation}\label{new:integer}
 \mathcal Z_N:=\Delta_N\mathcal R_N\in\operatorname{Mat}_3(\Z)
 \quad\hbox{for every }N\ge0.
\end{equation}
This is an explicit all-$N$ integrality theorem, not a definition by reduced
numerical denominators.

The same floor-chamber calculation as before gives a particularly simple
answer. Away from the small and exceptional primes already treated there,
if $e_3$ denotes its exponent profile, then the new profile is
\begin{equation}\label{new:profile}
 e_{34}(x)=e_3(x)+\mathbf1_{x>1/19}\qquad(x=N/p>0).
\end{equation}
On $x>1$ this means that the periodic profile is increased by one.
This is a finite exact claim: \texttt{verify\_arithmetic.py} checks all $270$
Farey-$29$ chambers, both initial and periodic parts, including every
integer-offset state at every floor wall. The extra harmonic cutoff uses
slope $10$, already in the existing wall set, so no new chambers are omitted.
The script uses only Python's integer and rational arithmetic.

For completeness, $E_{d4}(N,p)\le6\ell_p(29N+1)$: at each Legendre level,
$V_d\ge m_d-4$ and the loss is at most $m_d+2$.
Together with the first-functional bound, this proves uniformly that
$E_{34}\le\max(v_p(6),6\ell_p(29N+1))$. Small primes contribute $o(N)$,
and primes near chamber boundaries contribute $O(\log N)$, exactly as in
the previous proof. Applying the prime number theorem to \eqref{new:profile}
therefore gives
\begin{equation}\label{new:rate}
 \limsup_{N\to\infty}\frac{\log\Delta_N}{N}
 \le c_{34}=c_3+19
 =57.78399331966841174068\ldots<57.784<b.
\end{equation}
The added $19$ is the integral
$\int_{1/19}^{\infty}x^{-2}\,\dd x=19$.
There is also an independent integer-only upper certificate for the full
new integral in \texttt{verify\_arithmetic.py}; no fitted growth rate is used.

\section*{V. The irrationality proof}
\begin{theorem}\label{new:main}
The numbers $1,\zeta(3),\zeta(4)$ are linearly independent over $\Q$.
Consequently $A\zeta(3)+B\zeta(4)$ is irrational for every fixed
$(A,B)\in\Q^2\setminus\{(0,0)\}$.
\end{theorem}
\begin{proof}
Write the three rows of the integer matrix $\mathcal Z_N$ as
$\mathbf P_N^{(3)},\mathbf P_N^{(4)},\mathbf Q_N$.
Equations \eqref{new:errors} and \eqref{new:rate} show that
\[
 \|\zeta(3)\mathbf Q_N-\mathbf P_N^{(3)}\|
 +\|\zeta(4)\mathbf Q_N-\mathbf P_N^{(4)}\|
 \le e^{-(b-c_{34})N+o(N)}\longrightarrow0.
\]
Suppose $h_0+h_3\zeta(3)+h_4\zeta(4)=0$ with integers $h_i$ not all zero.
The integer row vector
\[
 \mathbf R_N=h_3\mathbf P_N^{(3)}+h_4\mathbf P_N^{(4)}+h_0\mathbf Q_N
 =h_3(\mathbf P_N^{(3)}-\zeta(3)\mathbf Q_N)
  +h_4(\mathbf P_N^{(4)}-\zeta(4)\mathbf Q_N)
\]
tends to zero componentwise. For all sufficiently large $N$, every component
has absolute value less than one, so every component is zero. But
\[
 \det\mathcal Z_N=36\Delta_N^3/\det V_N\ne0,
\]
which forces $(h_3,h_4,h_0)=(0,0,0)$, a contradiction.
If $A\zeta(3)+B\zeta(4)$ were rational, clearing its denominator and those
of $A,B$ would give just such a nonzero integer relation. This proves the
assertion.
\end{proof}

\noindent The invertibility argument is the required nonvanishing step.
It does not infer nonvanishing from an eigenvalue ratio or from a finite
numerical test. It also explains why using all three columns of the same
walk makes the proof simpler than trying to control the zeros of one column.

\section*{VI. Quantitative consequence and formalization interface}
For each fixed rational pair $(A,B)\ne(0,0)$, combine the integer numerator
rows with a fixed integer multiplier clearing $A,B$. Equations
\eqref{new:growth}--\eqref{new:rate} give numerator/denominator height at most
$e^{(a+c_{34})N+o(N)}$ and linear-form error at most
$e^{-(b-c_{34})N+o(N)}$. Reducing the fraction can only improve the resulting
approximation inequality. Thus the handoff's readouts have the common lower
bound
\[
 \liminf\delta_N(A,B)\ge\frac{b-c_{34}}{a+c_{34}}>0.015629
\]
along regular indices. The theorem ensures that their errors are nonzero.
This is a common lower bound, not a claim that all pairs have the same delta.

The proof separates into the following precise formalization obligations:
\begin{enumerate}
 \item the elementary integer-matrix lemma used in Part~V;
 \item the exact rational CMF product, determinant and concomitant identities;
 \item the two formal residue functionals and their boundary/transport identities;
 \item the explicit all-prime integer multiplier \eqref{new:deltaN};
 \item the finite rational floor certificate and its prime-number-theorem limit;
 \item the real spectral/contour upper estimates \eqref{new:errors}.
\end{enumerate}
The files verify finite algebraic and arithmetic identities with exact
arithmetic. They are not presented as Lean verification of the analytic
lemmas. No assertion is justified merely by sampled values of $N$.

\clearpage
\section*{VII. Complete underlying $\zeta(3)$ argument and certificates}
The remainder reproduces the underlying proof so that the preceding argument
has all its definitions, matrix identities, analytic estimates, and original
prime-profile table in this single document. In this part, $D_N$ denotes the
first-functional multiplier, whereas the simultaneous multiplier in
Parts~I--VI is $\Delta_N$.
\bigskip
\noindent The convergence rate is a spectral calculation. The substantive
arithmetic step is to use \emph{two formulas for the same rational numbers},
and take the better denominator bound separately at every prime. Each formula
alone gives a weaker estimate. Their combination proves the bound below.
The finite arithmetic certificate uses exact rational operations; a separate
integer-only certificate bounds its limiting integral.

\begin{theorem}\label{thm:main}
Consider the field $\mathrm{MeijerG}(5,2,5,5,1)$, with coordinate order
$(a_0,\ldots,a_4,b_0,\ldots,b_4)$, initial point
\[
 (1,1,0,0,2;0,0,0,0,0)
\]
and direction
\[
 (11,12,-16,-17,10;0,1,2,-1,-2).
\]
Write its forward walk as
\[
 W_N=\begin{pmatrix}V_N&*\\0&(-1)^N\end{pmatrix},\qquad
 T_N=V_N^{-T},\qquad
 U=\begin{pmatrix}0&-1&1\\0&0&2\end{pmatrix}.
\]
Define rational numbers $x_{Nj},y_{Nj}$ by
\[
 \begin{pmatrix}x_{N0}&x_{N1}&x_{N2}\\y_{N0}&y_{N1}&y_{N2}\end{pmatrix}=UT_N.
\]
The searched approximation is $j=0$, equivalently
\[
 \frac{(W_N^{-T})_{2,0}-(W_N^{-T})_{1,0}}{2(W_N^{-T})_{2,0}},
\]
with zero-based indices and initial readout rows $(0,-1,1,0)$ and $(0,0,2,0)$.
There are explicit positive integers $D_N$ such that all six numbers
$D_Nx_{Nj},D_Ny_{Nj}$ are integers, and
\[
 \log D_N\le (38.783993451+o(1))N.
\]
In particular, the same trajectory gives
\[
 \boxed{\delta>0.22521,\qquad \mu(\zeta(3))<5.4402.}
\]
Here the assertion about $\delta$ means that, for every fixed
$\delta_0<0.22521$, every defined readout $x_{Nj}/y_{Nj}$, reduced to
$p/q$, satisfies $|\z-p/q|\le |q|^{-1-\delta_0}$ for all sufficiently large $N$.
\end{theorem}

\section*{1. The analytic estimates}
Put $u=(0,-1,1)$ and $v=(0,0,2)$. The estimates needed below are
\begin{align}
 \|vT_N\|&\le \exp(aN+o(N)),\label{eq:qgrowth}\\
 \|\z vT_N-uT_N\|&\le \exp(-bN+o(N)),\label{eq:error}
\end{align}
where
\[
 a=53.28758332262020526735\ldots,\qquad
 b=59.51995592483447507214\ldots.
\]
These are the usual two spectral rates; they do not include denominator
clearing. For completeness, Section~6 identifies the relevant solution and
explains why the error estimate holds, including the two complex conjugate
characteristic roots.

Let $B(n)$ be the upper left $3\times3$ block of the single-step matrix,
so $V_N=B(0)\cdots B(N-1)$. With $H(n)=\operatorname{diag}(1,n,n^2)$,
\[
 H(n)B(n)H(n+1)^{-1}=A+O(n^{-1}).
\]
The rational matrix $A$ has one positive eigenvalue $\lambda_0=e^{-a}$,
and two nonreal conjugate eigenvalues of modulus $e^b$.
The exact characteristic polynomial is
\[
 \chi(X)=X^3-c_2X^2+c_1X-c_0,
\]
{\small
\begin{align*}
 c_2&=\frac{83220721427294106907399200607730930785765129529801897078934559}{1484577825190137274336280508756918272}\\
 c_1&=\frac{731334801243107466089356596505717853275074571608716447379201100334794667003154611901941573810757}{146466868948884247884332915718530654652342544}\\
 c_0&=\frac{76623012348334313797938377862237790794460627983667564193521662995163373351121226574589853696}{2130520759189220333584634002331004844478493702212809977860478179}\\
\end{align*}
}
Its discriminant is negative, and its unique real root is bracketed by
$7.202721273276511926282264010546996958\cdot10^{-24}$ and
$7.202721273276511926282264010546996959\cdot10^{-24}$.
The identities $a=-\log\lambda_0$ and $b=\tfrac12\log(c_0/\lambda_0)$
give the intervals below. The certificate bounds these logarithms using an
exact rational series with a rational remainder, without floating-point assumptions.
The matrix and independent algebraic reconstruction are included in the proof files. In particular,
\[
 53.28758332262<a<53.28758332263,\qquad
 59.51995592483<b<59.51995592484.
\]

\section*{2. The two exact formulas}
Throughout, $(t)_r=t(t+1)\cdots(t+r-1)$, with $(t)_0=1$.
First set
\[
 R_N(t)=\frac{(t-17N)_{15N}(t-16N)_{15N}}
 {(t)_{10N+2}(t+N)_{10N+1}(t+2N)_{10N+1}}.
\]
For $j=0,1,2$, expand by partial fractions the convergent sum
\[
 -\sum_{t=2N+1}^{\infty}\frac{\dd}{\dd t}\bigl((-t)^jR_N(t)\bigr)
 =A_{Nj}\zeta(4)+B_{Nj}\zeta(3)+C_{Nj}.
\]
The coefficients on the right are defined by partial fractions, so the
following identity is a coefficientwise identity over $\Q$:
\begin{equation}\label{eq:originalidentity}
 \mathcal C_N:=(A_{Nj},B_{Nj},C_{Nj})_{\text{rows}}
 =(-1)^N\mathcal C_0V_N,\qquad
 \mathcal C_0=\begin{pmatrix}3&0&0\\-2&-2&0\\1&1&1\end{pmatrix}.
\end{equation}
No independence assertion about $\zeta(3)$ and $\zeta(4)$ is used.
Also, exactly,
\begin{equation}\label{eq:det}
 \det V_N=F_N=
 \frac{\prod_{r\in\mathcal A}(rN)!}
 {\prod_{r\in\mathcal B}(rN)!\prod_{r\in\mathcal E}(rN+1)!},
\end{equation}
where repetitions in these lists are retained:
\begin{align*}
 \mathcal A&=(16,17,18,15,14,17,18,19,16,15),\\
 \mathcal B&=(11,10,9,12,13,12,11,10,13,14),\\
 \mathcal E&=(10,9,8,11,12).
\end{align*}
The identity follows either by the axis determinants or by checking the
rational identity $\det B(n)=F_{n+1}/F_n$ and $F_0=1$.

The second formula is considerably simpler. Put
\begin{align*}
 R_N^*(t)&=\frac{(t+N+1)_{11N}(t+1)_{11N}(t+2N)_{8N+1}}
 {(t+13N+1)_{15N+1}(t+14N+1)_{15N+1}},\\
 \kappa_N&=450N^2+55N+2,\qquad S_N^*(t)=(t+\kappa_N)R_N^*(t).
\end{align*}
Then, for every $N\ge0$,
\begin{equation}\label{eq:dualidentity}
 1-\sum_{t=-11N}^{\infty}(S_N^*)'(t)
       =(-1)^N(x_{N2}-\zeta(3)y_{N2}).
\end{equation}
The $1$ is necessary: $S_N^*(t)=1+O(t^{-2})$ at infinity. In particular,
at $N=0$ this is $1-2\zeta(3)$, exactly the initial readout.

Write
\[
 S_N^*(t)=1+\sum_k\left(\frac{d_{k1}}{t+k}+\frac{d_{k2}}{(t+k)^2}\right).
\]
The pole range is $13N+1\le k\le29N+1$, with double poles precisely on
$14N+1\le k\le28N+1$, and $\sum_kd_{k1}=0$. Thus
\begin{align*}
 (-1)^{N+1}y_{N2}&=2\sum_k d_{k2},\\
 (-1)^Nx_{N2}&=1-\sum_k\left(d_{k1}H_{k-11N-1}^{(2)}
                         +2d_{k2}H_{k-11N-1}^{(3)}\right).
\end{align*}
The largest harmonic index is $18N$. These are finite formulas for the
actual CMF coefficients, not a second, unrelated approximation sequence.

For use with all three columns, define the integer matrix
\[
 J_N=\begin{pmatrix}
 4N^2(13725N^2+2389N+103)&-6N(275N^2+150N+11)&-450N^2-43N+2&3\\
 N(5N+1)(330N+23)&N(450N+43)&-2&0\\
 -(18N+1)(25N+1)&1&0&0\\
 0&0&0&0
 \end{pmatrix}.
\]
Each readout is given by the same formal residue functional applied to
$Q_{Nj}(t)R_N^*(t)$, where
\[
 Q_{Nj}(t)=\sum_{r=0}^3(J_N)_{jr}(-t-12N-1)^r\in\Z[t].
\]
Its polynomial part has integer coefficients and degree at most two. Its
contribution therefore has denominator dividing $6$; the proper part has
the same double poles and harmonic indices as above. Consequently the
denominator estimate derived from this formula applies to all three columns.
Section~7 proves the coefficient identities and explains the polynomial term.

\section*{3. Explicit all-$N$ integer multipliers}
Here are fully explicit bounds, including prime powers and small primes.
For a prime $p$, let
\[
 f_p(m)=v_p(m!)=\sum_{r\ge1}\left\lfloor\frac m{p^r}\right\rfloor,
 \qquad \ell_p(m)=v_p(\lcm(1,\ldots,m)),\qquad \ell_p(0)=0.
\]
For an integer interval $I=[a,b]$ and an integer $k$, define
\[
 G_p(k;[a,b])=\begin{cases}
 f_p(k-a)+f_p(b-k),&a\le k\le b,\\
 f_p(b-k)-f_p(a-k-1),&k<a,\\
 f_p(k-a)-f_p(k-b-1),&k>b.
 \end{cases}
\]
This is exactly the valuation of $\prod_{r=a,\ r\ne k}^b(r-k)$.

For the original kernel use
\begin{align*}
 \mathcal I_o&=([0,10N+1],[N,11N],[2N,12N]),\\
 m_o(k)&=\#\{I\in\mathcal I_o:k\in I\},\\
 V_o(k)&=f_p(k+17N)+f_p(k+16N)-f_p(k+2N)-f_p(k+N)
              -\sum_{I\in\mathcal I_o}G_p(k;I),\\
 A_p&=\min_{2N\le k\le10N+1}V_o(k),\\
 C_p&=\min_{0\le k\le12N}
 \{V_o(k)-(m_o(k)-1)\ell_p(k+17N)-2\ell_p(k+2N)\},\\
 E_o(N,p)&=\max\{0,v_p(F_N)-A_p-C_p\}\qquad(p>3).
\end{align*}
For the dual kernel use
\begin{align*}
 \mathcal I_d&=([13N+1,28N+1],[14N+1,29N+1]),\\
 m_d(k)&=\#\{I\in\mathcal I_d:k\in I\},\\
 V_d(k)&=f_p(k-N-1)+f_p(k-1)+f_p(k-2N)\\
 &\quad-f_p(k-12N-1)-f_p(k-11N-1)-f_p(k-10N-1)
             -\sum_{I\in\mathcal I_d}G_p(k;I),\\
 h(k)&=k-11N-1,\qquad L(k)=\max\{k-1,29N+1-k\},\\
 E_d(N,p)&=\max\left\{v_p(6),0,
 \max_{13N+1\le k\le29N+1}
 \bigl((m_d(k)-1)\ell_p(L(k))+2\ell_p(h(k))-V_d(k)\bigr)\right\}.
\end{align*}
The second term in $L(k)$ must be retained: denominator differences can be
larger than $k-1$ near the left end of the double-pole interval.

For $N\ge1$, the promised explicit integer is
\begin{equation}\label{eq:DN}
 D_N=2^{E_d(N,2)}3^{E_d(N,3)}
       \prod_{3<p\le29N+1}p^{\min\{E_o(N,p),E_d(N,p)\}},
 \qquad D_0=1.
\end{equation}

\begin{lemma}\label{lem:integer}
For every integer $N\ge0$ and $j=0,1,2$, the numbers $D_Nx_{Nj}$ and
$D_Ny_{Nj}$ are integers.
\end{lemma}
\begin{proof}
At a pole $-k$ of order $m$, the top partial-fraction coefficient has
valuation $V_o(k)$ or $V_d(k)$ respectively. Each Taylor derivative can
lose at most $\ell_p(M)$, where $M$ bounds the absolute values of all
nonzero integer factors at that pole. Thus the coefficient of
$(t+k)^{-r}$ has valuation at least $V-(m-r)\ell_p(M)$.
Multiplication by an integer polynomial does not weaken this statement.
For the original kernel the possible factor $1/2$ from the second Taylor
coefficient is harmless because this bound is used only for $p>3$.

The rational coefficient contributed by that pole also includes
$H_h^{(r+1)}$, with valuation at least $-(r+1)\ell_p(h)$. Since
$\ell_p(M)\ge\ell_p(h)$, its valuation is at least
\[
 V-(m-1)\ell_p(M)-2\ell_p(h).
\]
This proves the bound $C_p$ for every entry of the rational row of
$\mathcal C_N$. The $\zeta(3)$ row satisfies the same weaker bound, and
the $\zeta(4)$ row has valuation at least $A_p$.
Now
\[
 u\mathcal C_0^T=(0,2,0),\qquad v\mathcal C_0^T=(0,0,2),\qquad
 \det\mathcal C_N=-6(-1)^NF_N.
\]
Consequently the cofactors expressing $x_{Nj}$ use the $\zeta(4)$ and
rational rows, and those expressing $y_{Nj}$ use the $\zeta(4)$ and
$\zeta(3)$ rows. This gives the denominator bound $E_o(N,p)$.

For the dual formula, the same pole argument, now with $m\le2$, proves
the bound $E_d(N,p)$ for the proper part at every prime. The integer
polynomial part has degree at most two, so its Bernoulli evaluation has
denominator dividing $6$, already included in $E_d$.
Both formulas describe the same rational numbers, so at each prime we
may take the minimum of the two bounds. No prime larger than $29N+1$
occurs. This proves the lemma; $N=0$ follows directly from $U$.
\end{proof}

\section*{4. The denominator growth calculation}
We explain why this is a finite floor calculation, rather than a finite-depth
experiment. For $p>\sqrt{30N+10}$ all factorial valuations in Section~3
equal $\lfloor m/p\rfloor$, and every $\ell_p(m)$ equals the indicator
of $m\ge p$. Put $x=N/p$ and $y=k/p$.

Discard temporarily the primes dividing
\[
 \prod_{d=1}^{29}\prod_{r=-10}^{10}(dN+r).
\]
For large $N$ this is a nonzero integer of polynomial growth. The logarithmic
contribution of these primes is $O(\log N)$: the exponents at the remaining
large primes are uniformly bounded. Small primes contribute $o(N)$, because
the explicit uniform bound is
\[
 E_d(N,p)\le\max\{v_p(6),5\ell_p(29N+1)\}.
\]
At each Legendre level the dual top-coefficient valuation is at least
$m_d(k)-4$: in the simple-pole case there are four positive and six negative
factorial floors with zero net linear argument, and in the double-pole case
three positive and seven negative floors with net linear argument one.
The loss is at most $m_d(k)+1$ per level. Thus small primes contribute
$O(\sqrt N\log N)$, and the exceptional large primes have bounded exponents.

Every possible jump of a factorial floor, pole membership, or loss indicator
as $k$ varies has the form $k=cN+jp+s$, with $|s|\le2$. The relevant sets
of $c$ are
\begin{align*}
 \mathcal S_o&=\{-17,-16,-2,-1,0,1,2,10,11,12\},\\
 \mathcal S_d&=\{0,1,2,10,11,12,13,14,28,29\}.
\end{align*}
The order of two distinct jumps can change only when
$(c-c')x$ is an integer. Since $|c-c'|\le29$, the Farey fractions of order
$29$ partition $(0,1)$ into all the required chambers. The discarded primes
cover the bounded integer offsets near chamber boundaries.

On each chamber it suffices to inspect the jumps and the constant regions
between them. The certificate also inspects offsets $-3,\ldots,3$ at every
jump using exact rational numbers with a formal positive infinitesimal.
This covers isolated endpoint states, not just open intervals.
There are $270$ chambers. Taking the minimum of the two resulting bounds
gives a step function $e_0(x)$ for $0<x<1$.

For $x>1$ the bound is periodic in $x$; denote it by $e_1(\{x\})$.
Indeed each pole-order sector then contains a full period of residues
$k\bmod p$, both loss indicators are one, and the factorial slope sums
balance. The factorial determinant also has balanced slope sum
$165-115-50=0$. The exact chamber calculation gives the table in the
appendix. This calculation uses only floors of rational numbers, comparisons,
and integer maxima and minima.

The prime number theorem now gives
\begin{align}\label{eq:c}
 \limsup_{N\to\infty}\frac{\log D_N}{N}
 &\le c:=\int_0^1\frac{e_0(x)}{x^2}\,\dd x
       +\int_0^1e_1(x)\psi'(1+x)\,\dd x\\
 &=38.78399331966841174068336120201250063785\ldots.\nonumber
\end{align}
For example, a chamber $(l,r)$ in the first integral contributes
$e_0(1/l-1/r)$, and its contribution from the intervals $(k+l,k+r)$,
$k\ge1$, is
$e_1\sum_{k\ge1}((k+l)^{-1}-(k+r)^{-1})$.
One first applies the prime number theorem on finitely many such intervals;
the remaining $x$-tail is $O(1/K)$ uniformly and can then be discarded.

There is no need to trust a floating-point digamma evaluation. The first
integral is exactly $9967051/272272$. Sum the rational interval terms in
the second integral for $k=1,\ldots,4096$, rounding each upward to a multiple
of $10^{-16}$. Its remaining tail is at most $m/4096$, where
\[
 m=\int_0^1e_1(x)\,\dd x=\frac{1313791}{604656}.
\]
This purely integer calculation gives
\begin{equation}\label{eq:cup}
 c\le38.7839934501989356<38.784.
\end{equation}
Together with Lemma~\ref{lem:integer}, this proves the arithmetic assertion
in Theorem~\ref{thm:main}.

This is also the desired GCD estimate. If $d_N$ is any integer clearing the
original rational readouts and $p_{Nj}^{\mathrm{raw}}=d_Nx_{Nj}$,
$q_{Nj}^{\mathrm{raw}}=d_Ny_{Nj}$, their least common denominator $L_{Nj}$
divides both $d_N$ and $D_N$. Therefore
\[
 \gcd(p_{Nj}^{\mathrm{raw}},q_{Nj}^{\mathrm{raw}})
 \ge\frac{d_N}{L_{Nj}}\ge\frac{d_N}{D_N},\qquad
 \log\gcd\ge\log d_N-cN-o(N).
\]

\section*{5. Passing from integer linear forms to $\mu$}
Put $P_{Nj}=D_Nx_{Nj}$ and $Q_{Nj}=D_Ny_{Nj}$. They are integers, and
\[
 |Q_{Nj}|\le e^{(a+c)N+o(N)},\qquad
 |Q_{Nj}\z-P_{Nj}|\le e^{-(b-c)N+o(N)}.
\]
The $2\times3$ matrix with these three pairs as columns has rank two,
because $U$ has rank two and $T_N$ is invertible. Thus for any rational
$h/k$ with $k>0$, at least one column satisfies
$kP_{Nj}-hQ_{Nj}\ne0$.
Choose $N$ of size $\log k/(b-c-\varepsilon)$ so that all three linear
forms have absolute value at most $1/(2k)$. For that column,
\[
 1\le |kP_{Nj}-hQ_{Nj}|
 \le\tfrac12+|Q_{Nj}|\,|k\z-h|.
\]
It follows, on letting $\varepsilon$ decrease to zero, that
\begin{equation}\label{eq:mu}
 \mu(\zeta(3))\le1+\frac{a+c}{b-c}
       =\frac{a+b}{b-c}.
\end{equation}
No lower bound on an oscillating error term is required.
Likewise the approximation exponent is at least
\[
 \delta_*:=\frac{b-c}{a+c}
       =0.22521567851203716390\ldots.
\]
Reducing a fraction can only improve the stated denominator-based
approximation inequality. With the rigorous upper bound
\eqref{eq:cup} and the spectral intervals in Section~1, the strict rounded
conclusions are $\delta>0.22521$ and $\mu(\zeta(3))<5.4402$.

\section*{6. Why the analytic estimates apply to this exact solution}
Let $\mathbf A_N,\mathbf B_N,\mathbf C_N$ be the three rows of
$\mathcal C_N$, and put
$\mathbf S_N=\zeta(4)\mathbf A_N+\zeta(3)\mathbf B_N+\mathbf C_N$.
The exact contour representation is
\[
 S_{Nj}=\frac1{2\pi i}\int_{2N+1/2-i\infty}^{2N+1/2+i\infty}
 \left(\frac\pi{\sin\pi t}\right)^2(-t)^jR_N(t)\,\dd t.
\]
Closing to the right gives the derivative sum in Section~2, with its displayed
minus sign. On this contour write $t=2N+\tfrac12+iy$. Directly from the factors,
\[
 |R_N(t)|\le \frac1{48N^4}
 \exp\left((15\log210-10\log24)N+\frac{29}{14}|y|\right).
\]
The factor $|\pi/\sin(\pi t)|^2$ decays as $e^{-2\pi|y|}$.
Consequently all three components of $\mathbf S_N$ are bounded by a
polynomial in $N$ times $e^{c_0N}$, where
\[
 c_0=15\log210-10\log24=48.4260746572\ldots<b.
\]
This excludes the two growing characteristic modes.

Here is an elementary justification even though these two roots are
nonreal. In a fixed real basis the limiting matrix is a scalar of modulus
$e^{-a}$ on one coordinate and $e^b$ times a rotation on the other two.
For sufficiently small perturbations, the cone in which the latter
two-coordinate norm is at least the first-coordinate modulus is forward
invariant and expands at any prescribed rate slightly less than $e^b$.
Our bound $e^{c_0N}$ prevents entry into that cone. Outside it, the first
coordinate bounds the whole vector and contracts at a rate arbitrarily
close to $e^{-a}$. Polynomial gauges do not change exponential rates. Thus
\[
 \|\mathbf S_N\|\le e^{-aN+o(N)}.
\]
Moreover $\|\mathbf A_N\|\le e^{bN+o(N)}$ and
$\det V_N=e^{(2b-a)N+o(N)}$. The exact cross-product identity is
\[
 (\z v-u)V_N^{-T}
       =-\frac{\mathbf A_N\times\mathbf S_N}{3\det V_N}.
\]
It proves \eqref{eq:error}; the inverse-transpose product bound proves
\eqref{eq:qgrowth}. These arguments give the required upper bounds without
any assumption about the phase of the complex modes.

\section*{7. Coefficient identities and the boundary constant}
The following elementary formalism makes the all-$N$ claims independent of
numerical recognition of zeta values. If all poles of $f$ are integers at
most $K$, write its polynomial and partial-fraction decomposition and define
\[
 \mathcal L_K\!\left[\sum_m p_mt^m+
               \sum_{h,r}\frac{c_{hr}}{(t-h)^r}\right]
 =\sum_m p_m B_m(K+1)+
       \sum_{h,r}r c_{hr}(Z_{r+1}-H_{K-h}^{(r+1)}).
\]
Here $Z_r$ are independent formal symbols, $B_1(x)=x-\tfrac12$, and
$B_m(x)$ are Bernoulli polynomials. For a proper rational function this
is the formal value of $-\sum_{t>K}f'(t)$. Its exact shift identities are
\[
 \mathcal L_K[f(t)-f(t-1)]=f'(K),\qquad
 \mathcal L_K[f]-\mathcal L_{K-1}[f]=f'(K),
\]
when the indicated functions have no pole at the boundary.
For proper parts these follow by telescoping harmonic sums; for polynomials
they are the identity $B_m(x+1)-B_m(x)=mx^{m-1}$.

After reflecting the first two upper gamma factors, the Meijer kernel is
$(\pi/\sin\pi t)^2r_{a,b}(t)$, with rational $r_{a,b}$. It satisfies
\[
 \frac{r(t)}{r(t-1)}=
 \frac{\prod_j(t-1+b_j)}{\prod_i(t-1+a_i)}.
\]
With $D(\theta)=\prod_i(\theta-a_i+1)-\prod_j(\theta-b_j)$ and
$g(t)=\prod_j(t+b_j)r(t)$, this gives
\[
 D(-t)r(t)=g(t)-g(t-1).
\]
Whenever $g'(K)=0$, the four formal moments of $r,-tr,t^2r,-t^3r$
therefore satisfy exactly the fourth-order companion relation used by the
CMF. The axis transformations follow directly from $\Gamma(t+1)=t\Gamma(t)$.

For the original parameters, first lower $a_2$ by $16N$ and $a_3$ by $17N$.
The resulting double zeros at $1,\ldots,2N$ allow the cutoff to move from
$0$ to $2N$. Then raise the three positive upper parameters, raise the
positive lower parameters, and finally lower the two negative lower
parameters. Along this path the leading companion coefficient
$5-\sum a_i+\sum b_j$ stays positive; every axis matrix is invertible.
The boundary factor in $g$ has a double zero throughout. This proves
\eqref{eq:originalidentity} by exact transport from $N=0$.

For the dual parameters $a'=1-a$, $b'=-b$, the initial rational kernel is
$(s-1)/s^2$ and its four formal moments are
\[
 g_0=(Z_2-2Z_3,\ Z_2-1,\ -\tfrac12,\ \tfrac13).
\]
First lower the three nonpositive upper parameters. Triple zeros at
$1,\ldots,N$ allow the cutoff to move to $K=N$. Then raise the two positive
upper parameters, raise the positive lower parameters, and lower the
negative lower parameters last. The companion coefficient remains positive
and all shifts are nonsingular; $g'(N)=0$ throughout. Thus these formal
moments have the exact dual CMF transport, denoted $D_N^{\mathrm{dual}}$.

The polynomial matrix $J_N$ in Section~2 is the Lagrange concomitant of
$z\prod_i(\theta+1-a_i)-\prod_j(\theta-b_j)$ and its formal adjoint. It obeys
the exact identity
\[
 W_NJ_N(D_N^{\mathrm{dual}})^T=J_0.
\]
This can also be verified as a polynomial identity for a generic elementary
upper shift and lower shift, then multiplied along the path. Its exact
symbolic certificate is included with the proof files. The shift identities
are first multiplied at generic parameters and only the complete transports
are then specialized. The original and dual nonsingular paths described
above show that both complete transports have regular specializations;
a common path need not have nonsingular intermediate charts on both sides.
Set $w_p=(-1,0,1/2,-1/3)$ and $w_q=(-2,0,0,0)$. Then
\[
 g_0=-w_p+Z_3w_q+(Z_2-1)(1,1,0,0),
\]
and $(1,1,0,0)$ is in the relevant kernel of $J_0$. Pairing and transport
therefore give $g_NJ_N[0\!:\!3,:]^T=Z_3vT_N-uT_N$ coefficientwise.

At the endpoint substitute $s=t+12N+1$. The kernel becomes
$(-1)^NR_N^*(t)$, and the first summation index becomes $t=-11N$.
The third row of $J_N$ multiplies the kernel by $-(t+\kappa_N)$.
Its polynomial part is $-1$, whose formal value is $-1$.
This proves \eqref{eq:dualidentity}, including its sign, normalization,
lower limit and boundary constant. The other two rows give the integer
polynomials $Q_{Nj}$ used above. This proof is entirely coefficientwise.

\appendix
\section*{Appendix: the exact prime-exponent profile}
In the following table $(l,r)$ is an open interval. The first exponent is
used for $0<N/p<1$, and the second for the fractional part of $N/p>1$.
Boundary primes have already been accounted for separately.
\begin{longtable}{rrrr}
\toprule $l$&$r$&$e_0$&$e_1$\\\midrule\endhead
$0$&$\frac{1}{18}$&0&3\\
$\frac{1}{18}$&$\frac{1}{17}$&1&3\\
$\frac{1}{17}$&$\frac{1}{16}$&2&3\\
$\frac{1}{16}$&$\frac{1}{15}$&3&3\\
$\frac{1}{15}$&$\frac{1}{13}$&4&4\\
$\frac{1}{13}$&$\frac{1}{12}$&3&3\\
$\frac{1}{12}$&$\frac{1}{11}$&2&2\\
$\frac{1}{11}$&$\frac{1}{10}$&1&1\\
$\frac{1}{10}$&$\frac{2}{19}$&0&0\\
$\frac{2}{19}$&$\frac{1}{8}$&1&1\\
$\frac{1}{8}$&$\frac{2}{15}$&2&2\\
$\frac{2}{15}$&$\frac{2}{13}$&3&3\\
$\frac{2}{13}$&$\frac{3}{17}$&2&2\\
$\frac{3}{17}$&$\frac{2}{11}$&3&3\\
$\frac{2}{11}$&$\frac{3}{13}$&2&2\\
$\frac{3}{13}$&$\frac{4}{17}$&1&1\\
$\frac{4}{17}$&$\frac{3}{11}$&2&2\\
$\frac{3}{11}$&$\frac{5}{18}$&1&1\\
$\frac{5}{18}$&$\frac{5}{17}$&2&2\\
$\frac{5}{17}$&$\frac{3}{10}$&3&3\\
$\frac{3}{10}$&$\frac{4}{11}$&2&2\\
$\frac{4}{11}$&$\frac{7}{19}$&1&1\\
$\frac{7}{19}$&$\frac{5}{13}$&2&2\\
$\frac{5}{13}$&$\frac{7}{18}$&1&1\\
$\frac{7}{18}$&$\frac{5}{12}$&2&2\\
$\frac{5}{12}$&$\frac{8}{19}$&1&1\\
$\frac{8}{19}$&$\frac{5}{11}$&2&2\\
$\frac{5}{11}$&$\frac{8}{17}$&1&1\\
$\frac{8}{17}$&$\frac{9}{19}$&2&2\\
$\frac{9}{19}$&$\frac{1}{2}$&3&3\\
$\frac{1}{2}$&$\frac{8}{15}$&2&2\\
$\frac{8}{15}$&$\frac{7}{13}$&3&3\\
$\frac{7}{13}$&$\frac{6}{11}$&2&2\\
$\frac{6}{11}$&$\frac{9}{16}$&1&1\\
$\frac{9}{16}$&$\frac{7}{12}$&2&2\\
$\frac{7}{12}$&$\frac{10}{17}$&1&1\\
$\frac{10}{17}$&$\frac{7}{11}$&2&2\\
$\frac{7}{11}$&$\frac{11}{17}$&1&1\\
$\frac{11}{17}$&$\frac{11}{16}$&2&2\\
$\frac{11}{16}$&$\frac{9}{13}$&3&3\\
$\frac{9}{13}$&$\frac{8}{11}$&2&2\\
$\frac{8}{11}$&$\frac{11}{15}$&1&1\\
$\frac{11}{15}$&$\frac{14}{19}$&2&2\\
$\frac{14}{19}$&$\frac{3}{4}$&3&3\\
$\frac{3}{4}$&$\frac{13}{16}$&2&2\\
$\frac{13}{16}$&$\frac{9}{11}$&3&3\\
$\frac{9}{11}$&$\frac{16}{19}$&2&2\\
$\frac{16}{19}$&$\frac{7}{8}$&3&3\\
$\frac{7}{8}$&$\frac{15}{17}$&2&2\\
$\frac{15}{17}$&$\frac{9}{10}$&3&3\\
$\frac{9}{10}$&$\frac{10}{11}$&2&2\\
$\frac{10}{11}$&$\frac{11}{12}$&1&1\\
$\frac{11}{12}$&$\frac{15}{16}$&0&0\\
$\frac{15}{16}$&$\frac{17}{18}$&2&2\\
$\frac{17}{18}$&$\frac{18}{19}$&3&3\\
$\frac{18}{19}$&$1$&4&4\\
\bottomrule
\end{longtable}

\noindent\textbf{Certificate scope.} The denominator formula and the
coefficient identities are all-$N$ arguments. The $270$ floor chambers and
their endpoint states are checked with exact arithmetic. The upper bound
$c<38.784$ also has an integer-only certificate. Numerical trajectory
checks were confined to depth $500$; none is used as a substitute for an
all-$N$ identity or for the arithmetic growth argument.
\end{document}
